\documentclass[12pt]{article}
\usepackage{amsmath, amssymb, amsthm}
\usepackage{mathtools}
\usepackage{hyperref}
\usepackage{listings}
\usepackage{xcolor}

\lstdefinelanguage{lean}{
  keywords={import, namespace, open, variable, def, lemma, theorem, structure, inductive, deriving, section, end},
  keywordstyle=\color{blue}\bfseries,
  ndkeywords={sorry, by, exact},
  ndkeywordstyle=\color{red}\bfseries,
  comment=[l]{--},
  commentstyle=\color{gray}\ttfamily,
  stringstyle=\color{green}\ttfamily,
  morestring=[b]",
  basicstyle=\ttfamily\small,
  breaklines=true
}

\title{BOSONIZE-LEAN: Phase 1 Formalization Blueprint}
\author{Yellow Team Architect}
\date{\today}

\begin{document}
\maketitle

\section{Introduction and Architecture}
Based on the Red Team Phase 1 critique, the monolithic proofs and trivial stubs (\texttt{0}, \texttt{id}, \texttt{bottom}) are strictly banned. This document strictly mirrors the required Lean 4 stubs in \texttt{Bosonize/Stubs/Phase1\_Stubs.lean}. 

\section{1.1 Lattice and Index Types}
We define the spatial lattice as a periodic ring and the momentum band as a centered interval.
\begin{lstlisting}[language=lean]
def Lambda (L : Nat) [NeZero L] := ZMod L
def LambdaDual (L : Nat) := {n : Int // - (L : Int) < 2 * n /\ 2 * n <= (L : Int)}
lemma card_LambdaDual (hL : Even L) : Fintype.card (LambdaDual L) = L := sorry
inductive Chirality | R | L
\end{lstlisting}

\section{1.2 Umbral Calculus Core}
The continuous derivatives are replaced entirely by finite differences on the integers.
\begin{lstlisting}[language=lean]
def shiftE (f : Int -> R) : Int -> R := fun x => f (x + 1)
def fwdDiff (f : Int -> R) : Int -> R := fun x => f (x + 1) - f x
def bwdDiff (f : Int -> R) : Int -> R := fun x => f x - f (x - 1)

lemma discrete_leibniz (f g : Int -> R) (x : Int) :
  fwdDiff (f * g) x = fwdDiff f x * g x + shiftE f x * fwdDiff g x := sorry

lemma sum_fwdDiff_eq_bounds (f : Int -> R) (a b : Int) (hab : a <= b) :
  \sum x \in Ico a b, fwdDiff f x = f b - f a := sorry

lemma sum_by_parts_Lambda (f g : Lambda L -> R) :
  \sum x : Lambda L, fwdDiff f x * g x = - \sum x : Lambda L, f x * bwdDiff g x := sorry

def umbralBeta (f : Int -> R) : Int -> R := fun x => (x : R) * f (x - 1)
lemma umbral_heisenberg (f : Int -> R) (x : Int) :
  fwdDiff (umbralBeta f) x - umbralBeta (fwdDiff f) x = f x := sorry
\end{lstlisting}

\section{1.3 Finite Fourier}
Fourier modes on the periodic lattice.
\begin{lstlisting}[language=lean]
def plane_wave (k : LambdaDual L) (x : Lambda L) : Complex :=
  \zeta ^ ((k : Int) * (x.val : Int))

lemma fwdDiff_plane_wave (k : LambdaDual L) (x : Lambda L) :
  fwdDiff (fun z => \zeta ^ ((k : Int) * z)) x.val = (\zeta ^ (k : Int) - 1) * plane_wave \zeta k x := sorry

lemma plane_wave_ortho (k k' : LambdaDual L) :
  \sum x : Lambda L, (plane_wave \zeta k x) * conj (plane_wave \zeta k' x) = if k = k' then (L : Complex) else 0 := sorry
\end{lstlisting}

\section{1.4 CAR Representation}
The explicit power-set implementation of the fermionic Fock space using the exact phase counting function.
\begin{lstlisting}[language=lean]
structure CARRep (i : Type) (V : Type) where
  c : i -> Module.End Complex V
  cdag : i -> Module.End Complex V
  anticomm_c_cdag : \forall i j, c i * cdag j + cdag j * c i = if i = j then 1 else 0
  anticomm_c_c : \forall i j, c i * c j + c j * c i = 0

def FockSpace (i : Type) := Finset i -> Complex

def sign_func (i : i) (S : Finset i) : Complex :=
  (-1) ^ (S.filter (fun k => k < i)).card

def concrete_CAR : CARRep i (FockSpace i) := sorry
lemma dim_fock_space : Module.finrank Complex (FockSpace i) = 2 ^ Fintype.card i := sorry
\end{lstlisting}

\section{1.5 Lattice AQFT Net}
The local subalgebras. The Red Team flagged that $\bot$ was used previously. We strictly require \texttt{Algebra.adjoin}.
\begin{lstlisting}[language=lean]
def Net (I : Finset (Lambda L)) : Subalgebra Complex (Module.End Complex (FockSpace (Chirality x Lambda L))) :=
  Algebra.adjoin Complex (sorry) 

lemma net_isotony (h : I \subseteq J) : Net I \le Net J := sorry
lemma net_twisted_locality (h_disj : Disjoint I J) (A : Net I) (B : Net J) : sorry
\end{lstlisting}

\section{1.6 \& 1.7 Vacuum and Energy Budget}
\begin{lstlisting}[language=lean]
def vacuum_state : FockSpace (Chirality x LambdaDual L) := sorry
def energy_func (S : Finset (LambdaDual L)) : Int := sorry
lemma energy_nonneg (S : Finset (LambdaDual L)) : 0 \le energy_func S := sorry
def BudgetSpace (K Nmax : Nat) : Submodule Complex (FockSpace (Chirality x LambdaDual L)) := sorry
\end{lstlisting}

\section{1.8 Umbral Boson Fock Layer}
\begin{lstlisting}[language=lean]
def BosonFock := MvPolynomial (Fin Q) Complex

def a_dag (m : Fin Q) : Module.End Complex BosonFock :=
  LinearMap.mulLeft Complex (MvPolynomial.X m)

def a_ann (m : Fin Q) : Module.End Complex BosonFock :=
  MvPolynomial.pderiv m

lemma boson_ccr (m n : Fin Q) :
  a_ann m * a_dag n - a_dag n * a_ann m = if m = n then 1 else 0 := sorry

lemma nilpotent_on_budget (K : Nat) (p : BosonFock) (hp : MvPolynomial.totalDegree p \le K) : 
  (sorry : BosonFock) = 0 := sorry
\end{lstlisting}

\end{document}
